\chapter{التناظر، والانحلال، وسلالم العزم الزاوي}
\providecommand{\ket}[1]{\left\lvert#1\right\rangle}

\section{التناظرات ومستويات الطاقة المنحلّة}

التناظر هو عملية لا تغيّر الوصف الفيزيائي لمنظومة ما. فهو لا يغيّر الهاملتوني، ولا معادلة Schr\"odinger، ولا الطاقات الممكنة. وإذا كان مؤثر وحدوي \(S\) يمثّل هذه العملية، فإن الثبات يعني
\begin{equation}
S H S^\dagger=H,
\end{equation}
أو بصورة مكافئة
\begin{equation}
[S,H]=0.
\end{equation}

الانسحاب مثال على ذلك. فإذا نقلت ذرة من موضع إلى آخر فلا تتغيّر طاقتها. وهذا صحيح لكل حالة ذرية، لا للحالة الأرضية وحدها. والدوران مثال آخر، وسيكون الموضوع الرئيسي هنا.

إن تبديل جسيمين متماثلين هو أيضًا تناظر. فتسمية إلكترونين بالإلكترون الأول والإلكترون الثاني، أو Harry وFred، لا تجعلهما مختلفين. ولا يمكن لتبديل التسميتين أن يغيّر خواص الذرة.

للبلورة نوع إضافي من تناظر الانسحاب. فإزاحة البلورة كلها عبر الفضاء انسحاب عادي. أمّا في البلورة اللانهائية، فإن الانسحاب بمقدار مسافة شبكية واحدة يعيد أيضًا الترتيب الفيزيائي نفسه. وبصورة مكافئة، يمكن إعادة وسم مواقع الشبكة بإزاحة فهرس كل موقع بمقدار واحد. وتفسد الحدود هذه العبارة الدقيقة في بلورة منتهية، لكنها تصبح تناظرًا جيدًا بعيدًا عن طرفي بلورة طويلة جدًا.

تجانس الفضاء---أي إن كل موضع يشبه أي موضع آخر---هو تناظر الانسحاب. وتماثل الفضاء في جميع الاتجاهات---أي إن كل اتجاه يشبه أي اتجاه آخر---هو تناظر الدوران.

يرتبط الانحلال بالتناظر، لكنه ليس الشيء نفسه. يكون مستوى الطاقة منحلاً عندما تكون هناك أكثر من حالة واحدة لها هذه الطاقة على وجه الدقة. وبالنسبة إلى هاملتوني عام ذي طيف متقطع، فإن المساواة التامة بين مستويين غير مرتبطين ستكون مصادفة لافتة. وحتى الاتفاق إلى جزء واحد من \(10^{15}\) ليس مساواةً تامة. وتشير حالات الانحلال التام عادةً إلى وجود تناظر.

يفضي التناظر أحيانًا إلى الانحلال، ولكن ليس دائمًا. ويبيّن المثال الأول معًا لماذا يمكن أن يحدث ذلك ولماذا قد لا يكفي تناظر واحد وحده.

\section{جسيم على دائرة}

تأمّل خرزة مقيّدة بالحركة على سلك دائري. يمكن وصف موضعها بالإحداثيين الديكارتيين \(x\) و\(y\)، لكن إحداثيًا زاويًا واحدًا يكفي:
\begin{equation}
\psi=\psi(\theta).
\end{equation}
وكثافة احتمال الموضع على الدائرة هي \(\psi^*(\theta)\psi(\theta)\).

يحوِّل دوران فعلي عكس عقارب الساعة بزاوية صغيرة \(\epsilon\) الدالة كما يلي:
\begin{equation}
\psi(\theta)\longrightarrow\psi(\theta-\epsilon).
\end{equation}
ولذلك
\begin{align}
\delta\psi(\theta)
&=\psi(\theta-\epsilon)-\psi(\theta)\\
&=-\epsilon\frac{d\psi}{d\theta}
+\mathcal O(\epsilon^2).
\end{align}

المؤثر
\begin{equation}
L_z=-i\hbar\frac{d}{d\theta}
\end{equation}
هيرميتي على الدوال الموجية الدورية. وهو النظير الزاوي لـ
\begin{equation}
p_x=-i\hbar\frac{d}{dx}.
\end{equation}
ووجود العامل \(i\) أساسي للهيرميتية. وبدلالة \(L_z\)، لدينا
\begin{equation}
\delta\psi
=-\frac{i\epsilon}{\hbar}L_z\psi
+\mathcal O(\epsilon^2),
\end{equation}
ومن ثم فإن مؤثر الدوران متناهي الصغر هو
\begin{equation}
R(\epsilon)
=1-\frac{i\epsilon}{\hbar}L_z
+\mathcal O(\epsilon^2).
\end{equation}
وهذا يحدد \(L_z\) بوصفه مولِّد الدورانات حول المحور العمودي على الدائرة.

\begin{classroomqa}
\lecturetimestamp{00:11:19}
\audiencequestion{ألا ينبغي أن تكون لإحدى العوامل التي تتضمن \(i\) إشارة معاكسة؟}
\lecturerresponse{يؤدي فحص الإشارات إلى العلاقة أعلاه. ومع الاصطلاح \(\psi(\theta)\to\psi(\theta-\epsilon)\) يكون التغيّر \(-\epsilon\,d\psi/d\theta\)، وهو يساوي \(-i\epsilon L_z\psi/\hbar\).}
\end{classroomqa}

لنوجد الآن الحالات الذاتية للعزم الزاوي:
\begin{equation}
L_z\psi=M\psi.
\end{equation}
وبالتعويض بالمؤثر التفاضلي نحصل على
\begin{equation}
-i\hbar\frac{d\psi}{d\theta}=M\psi,
\end{equation}
وحلولها من الشكل
\begin{equation}
\psi(\theta)=C\exp\left(\frac{iM\theta}{\hbar}\right).
\end{equation}

يوجد شرط إضافي لا يظهر في الحركة على خط لانهائي. فالنقطتان \(\theta=0\) و\(\theta=2\pi\) هما النقطة نفسها، ولذلك يجب على الدالة الموجية المدارية الأحادية القيمة أن تحقق
\begin{equation}
\psi(\theta+2\pi)=\psi(\theta).
\end{equation}
وبالنسبة إلى الدالة الذاتية أعلاه، يقتضي ذلك
\begin{equation}
\exp\left(\frac{2\pi iM}{\hbar}\right)=1.
\end{equation}
وعليه،
\begin{equation}
M=m\hbar,
\qquad
m\in\mathbb Z.
\end{equation}
ومن ثم فإن العزوم الزاوية المدارية الممكنة على الدائرة هي
\begin{equation}
\ldots,-2\hbar,-\hbar,0,\hbar,2\hbar,\ldots.
\end{equation}

هل يجب أن تكون للحالة ذات العزم الزاوي \(m\hbar\) الطاقة نفسها التي للحالة ذات العزم الزاوي \(-m\hbar\)؟ يصف \(m\) الموجب والسالب جهتي دوران متعاكستين. والسؤال هو ما إذا كان تناظر الدوران وحده يقتضي أن تكون لهاتين الحركتين طاقة متساوية.

\begin{classroomqa}
\lecturetimestamp{00:22:51}
\audiencequestion{هل يعني استخدام \(\psi(\theta-\epsilon)\) أن الجسيم يتحرك مع عقارب الساعة؟}
\lecturerresponse{لا. فهذا الاصطلاح يحدِّد فقط أيَّ اتجاه يُسمّى الدوران الموجب. وهو يعيّن إشارات الموضع الزاوي والعزم الزاوي، لكنه لا يحدِّد الاتجاه الفعلي الذي يتحرك فيه الجسيم.}
\end{classroomqa}

يوفّر مجال مغناطيسي عمودي على مستوى الدائرة مثالًا مضادًا. إذ يمكنه أن يزيح طاقات العزم الزاوي الموجب والسالب في اتجاهين متعاكسين، ومن ثم يفصل بينهما. وتعتمد إشارة هذا الانزياح على اصطلاحي الشحنة والاتجاه، لكن المهم هو حدوث هذا الانقسام.

يحافظ مثل هذا المجال مع ذلك على الدورانات حول المحور العمودي. ولذلك فهو يبيّن أن تناظر الدوران وحده لا يقتضي
\begin{equation}
E(m)=E(-m).
\end{equation}
وفي غياب هذا الانقسام، قد يشكّل \(m=1\) و\(m=-1\) زوجًا منحلاً، وكذلك قد يشكّل \(m=2\) و\(m=-2\) زوجًا منحلاً. أمّا الحالة \(m=0\) فليس لها قرين متميّز من هذا النوع.

\section{تناظر الانعكاس والزوج \texorpdfstring{\(m,-m\)}{m, -m}}

لنضف تناظر انعكاس متقطعًا. ولتدلّ \(\mathcal M\) على الانعكاس عبر قطر من أقطار الدائرة:
\begin{equation}
(\mathcal M\psi)(\theta)=\psi(-\theta).
\end{equation}
وهو يعكس جهة الدوران:
\begin{equation}
\mathcal M\psi_m=\psi_{-m}.
\end{equation}

إذا كان الانعكاس تناظرًا، فإن
\begin{equation}
[\mathcal M,H]=0.
\end{equation}
وافترض أن
\begin{equation}
H\ket{m}=E_m\ket{m}.
\end{equation}
فينتج من تطبيق \(\mathcal M\)
\begin{align}
H\mathcal M\ket{m}
&=\mathcal M H\ket{m}\\
&=E_m\mathcal M\ket{m}.
\end{align}
وبما أن \(\mathcal M\ket{m}\) هي الحالة ذات العزم الزاوي \(-m\hbar\)، فإن
\begin{equation}
E_{-m}=E_m.
\end{equation}
يحفظ تناظر الدوران العزم الزاوي؛ ويربط تناظر الانعكاس بين القيمتين المتعاكستين. ومعًا يقتضيان أن يكون الزوج منحلاً.

\begin{classroomqa}
\lecturetimestamp{00:29:54}
\audiencequestion{بما أن للنمطين \(m\) و\(-m\) كثافة الاحتمال نفسها، أفيدل ذلك وحده على أن طاقتيهما متساويتان؟}
\lecturerresponse{لا. فالدالتان الموجيتان \(e^{im\theta}\) و\(e^{-im\theta}\) مختلفتان على الرغم من أن مربعي مقداريهما متساويان. ومعادلة القيمة الذاتية للعزم الزاوي تتعلق بالدالة الموجية نفسها، لا بـ \(\psi^*\psi\) وحده. وتقتضي ميكانيكا الكم السعات، لا الاحتمالات وحدها.}
\end{classroomqa}

يحافظ المجال المغناطيسي على الدورانات حول محوره، لكنه يكسر تناظر الانعكاس المستخدم أعلاه.

\begin{classroomqa}
\lecturetimestamp{00:33:05}
\audiencequestion{هل يعكس تناظر الانعكاس المجال المغناطيسي؟}
\lecturerresponse{نعم. فالمجال المغناطيسي متجه محوري. وعند عكس مجال متجه إلى داخل الصفحة يصبح متجهًا إلى خارجها. وبصورة مكافئة، تحمل حلقة تيار وصورتها المرآتية تيارًا في جهتين متعاكستين، وتنتجان مجالين متعاكسين في الاتجاه.}
\end{classroomqa}

يمكن رؤية عدم التبادلية بين الدوران والانعكاس مباشرةً. فعلى دالة ذاتية للعزم الزاوي،
\begin{align}
\mathcal M L_z\psi_m
&=m\hbar\,\psi_{-m},\\
L_z\mathcal M\psi_m
&=-m\hbar\,\psi_{-m}.
\end{align}
ومن ثم
\begin{equation}
[\mathcal M,L_z]\psi_m
=2m\hbar\,\psi_{-m},
\end{equation}
وهذا غير صفري عندما يكون \(m\neq0\). فالدوران ثم الانعكاس ليس العملية نفسها كالانعكاس ثم الدوران.

يتبادل كلا المؤثرين مع \(H\) حين يكون التحويلان كلاهما تناظرين، لكنهما لا يتبادلان أحدهما مع الآخر. وهذا هو النمط المسؤول عن الانحلال.

\begin{classroomqa}
\lecturetimestamp{00:40:28}
\audiencequestion{هل يولِّد أي زوج غير متبادل من نوع Heisenberg، مثل \(x\) و\(p\)، حالات انحلال؟}
\lecturerresponse{ليس على العموم. فلا بد أيضًا أن يكون المؤثران تناظرين: أي يجب أن يتبادلا مع الهاملتوني. وعادةً لا يكون الموضع والزخم محفوظين معًا. ولو كان كلاهما يتبادل مع \(H\)، لأفادت عدم تبادليتهما وجود انحلال.}
\end{classroomqa}

\section{جبر المبدلات للتناظرات}

لنفترض أن \(A\) و\(B\) هما مولِّدان هيرميتيان لتناظرين متصلين:
\begin{equation}
[A,H]=0,
\qquad
[B,H]=0.
\end{equation}
ولا يقتضي شيء في هاتين المعادلتين أن يتبادل \(A\) و\(B\) أحدهما مع الآخر.

إن المبدِّل \([A,B]\) مضادّ للهيرميتية، لذا عرّف مؤثرًا هيرميتيًا
\begin{equation}
C=i[A,B].
\end{equation}
والعوامل العددية والإشارة الكلية لا أهمية لها هنا. والنقطة المهمة هي أن \(C\) يتبادل أيضًا مع الهاملتوني. وحقًّا،
\begin{equation}
[AB,H]
=A[B,H]+[A,H]B=0,
\end{equation}
وبالمثل
\begin{equation}
[BA,H]=0.
\end{equation}
ولذلك
\begin{equation}
[C,H]
=i\bigl([AB,H]-[BA,H]\bigr)=0.
\end{equation}

قد يكون المولِّد الجديد \(C\) مستقلًا عن \(A\) و\(B\)، وقد يكون مجرد تركيب خطي لمولِّدات معروفة من قبل. فإذا كان جديدًا فأضِفه إلى القائمة وشكِّل مبدِّلات أخرى. وقد يستمر تكرار هذه العملية إلى ما لا نهاية، لكنه كثيرًا ما ينغلق بعد عدد منتهٍ من المولِّدات: أي إن كل مبدِّل لاحق يكون تركيبًا خطيًا من تلك الموجودة سلفًا. وهذا النسق المغلق هو جبر مبدِّلات، أو جبر لي.

تشكّل تحويلات التناظر المنتهية زمرة. وتشكّل مولِّداتها متناهية الصغر جبر لي الموافق.

تكتسب حالات الانحلال أهمية لأنها تنظم الأطياف. فإذا كانت لحالتين الطاقة نفسها، فلا يمكن لإحداهما أن تضمحل إلى الأخرى بإصدار فوتون يحمل فرق الطاقة بينهما. وفي التحليل الطيفي الذري يحدِّد ترتيب المستويات وانحلالها الخطوط الطيفية الممكنة. وينطبق المنطق نفسه على الهادرونات المثارة وانتقالاتها.

\begin{classroomqa}
\lecturetimestamp{00:49:07}
\audiencequestion{هل هذه المبدلات نفسها عناصر في زمرة التناظر؟}
\lecturerresponse{المؤثرات التي يدور النقاش حولها هي مولدات. وتشكل المولدات ومبدلاتها جبر مبدلات، أو جبر لي. أما عمليات التناظر المنتهية المبنية من المولدات فهي عناصر الزمرة.}
\end{classroomqa}

إذا كانت جميع مولدات التناظر تتبادل مع بعضها بعضًا، سُمِّي التناظر أبيلِيًّا. وإذا كانت بعض مبدلاتها غير منعدمة، فهو غير أبيلِي.

\begin{classroomqa}
\lecturetimestamp{00:50:25}
\audiencequestion{ما التعريف الأساسي للمولد؟}
\lecturerresponse{خذ تحويل تناظر وحدويًا قريبًا جدًا من الهوية واكتبه على الصورة \(1-i\epsilon G/\hbar+\cdots\). والمؤثر الهيرميتي \(G\) هو مولده. ويؤدي تكرار التحويل المتناهي الصغر الموافق إلى بناء تحويل منتهٍ.}
\end{classroomqa}

\begin{classroomqa}
\lecturetimestamp{00:52:08}
\audiencequestion{ما المقصود بجبر مولدات التناظر؟}
\lecturerresponse{إنه مجموعة من المولدات يمكن جمعها وطرحها وضربها بأعداد عادية وأخذ مبدلاتها. وفي هذا الجبر تؤدي المبدلة دور الضرب، ويعني الانغلاق أن هذه العمليات لا تنتج مولدات خارج المجموعة.}
\end{classroomqa}

\section{الدورانات حول محاور مختلفة}

تمثل الدورانات أهم مثال على تناظر غير أبيلِي. فالدوران حول المحور \(x\) يتبعه دوران حول المحور \(y\) لا يساوي عمومًا إجراءهما بالترتيب المعكوس. وهذا الاختلاف بينهما يرتبط بدوره بدوران حول المحور الثالث.

\begin{classroomqa}
\lecturetimestamp{00:54:22}
\audiencequestion{هل تتبادل الدورانات المتناهية في الصغر رغم أن الدورانات المنتهية لا تتبادل؟}
\lecturerresponse{إنها تختلف عند المرتبة الثانية في زاويتي الدوران الصغيرتين. وهذا يكفي: فبعد فصل هذين المعلمَين الصغيرين، يكون للمولدات نفسها مبدلة غير منعدمة.}
\end{classroomqa}

قبل حساب المبدلات الكمومية، انظر إلى الصورة الكلاسيكية المناظرة. خذ مدارًا كوكبيًا دائريًا يكون فيه الزخم الزاوي عموديًا على مستواه. إن تدوير المدار كله حول المحور \(y\) يعطي مدارًا مسموحًا آخر له الطاقة نفسها تمامًا. غير أن متجه زخمه الزاوي صار الآن ذا مركبة \(x\) مختلفة.

وهكذا يولد الثبات الدوراني تهيئات لها الطاقة نفسها ومركبات مختلفة للزخم الزاوي. وكموميًا، يظهر هذا على هيئة عدم تبادلية بين مولدات الدوران حول محاور مختلفة.

\section{الزخم الزاوي في الإحداثيات الديكارتية}

لننظر أولًا في جسيم يتحرك في المستوى \(xy\)، من دون تقييده بدائرة. تحت دوران صغير عكس اتجاه عقارب الساعة،
\begin{equation}
\delta x=-\epsilon y,
\qquad
\delta y=\epsilon x.
\end{equation}
وبالنسبة إلى الحالة المُدارة فعليًا، تُقيَّم الدالة الموجية المتحولة عند الإحداثيات المدارة دورانًا عكسيًا:
\begin{equation}
\psi_\epsilon(x,y)
=\psi(x+\epsilon y,y-\epsilon x).
\end{equation}
ومن ثم،
\begin{align}
\delta\psi
&=\epsilon\left(
y\frac{\partial}{\partial x}
-x\frac{\partial}{\partial y}
\right)\psi
+\mathcal O(\epsilon^2)\\
&=-\frac{i\epsilon}{\hbar}
\left(xp_y-yp_x\right)\psi
+\mathcal O(\epsilon^2).
\end{align}
لذلك يكون مولد الدورانات حول المحور \(z\)
\begin{equation}
L_z=xp_y-yp_x.
\end{equation}
وهذه هي المركبة \(z\) للمتجه الكلاسيكي المألوف
\begin{equation}
\mathbf L=\mathbf r\times\mathbf p.
\end{equation}
فهو ليس مجرد الزخم الزاوي الميكانيكي؛ بل هو في ميكانيكا الكم المؤثر الذي يولد دورانات الدالة الموجية.\footnote{اختير اصطلاح الدوران الفعال ليتوافق مع التحويل السابق \(\psi(\theta)\mapsto\psi(\theta-\epsilon)\) ومع التعريف النهائي \(L_z=xp_y-yp_x\) بعد الالتباس الوسيط في الإشارة في الاشتقاق الشفهي.}

ويعطي التبديل الدوري \(x\to y\to z\to x\) المركبتين الأخريين.\footnote{تُستعاد صيغتا \(L_x\) و\(L_y\) بتطبيق التبديل الدوري \(x\to y\to z\to x\)، كما نُصَّ عليه صراحة في المحاضرة، على \(L_z=xp_y-yp_x\).}
\begin{align}
L_x&=yp_z-zp_y,\\
L_y&=zp_x-xp_z,\\
L_z&=xp_y-yp_x.
\end{align}

في نظام ثابت تحت الدوران، تتبادل هذه المؤثرات مع الهاملتوني:
\begin{equation}
[H,L_x]=[H,L_y]=[H,L_z]=0.
\end{equation}
ويمثل جسيم في مجال قوة مركزية المثال القياسي. وهذه المعادلات هي انحفاظ الزخم الزاوي بصيغته الكمومية.

\section{جبر لي للزخم الزاوي}

علاقات التبديل القانونية هي
\begin{align}
[x_i,x_j]&=0,\\
[p_i,p_j]&=0,\\
[x_i,p_j]&=i\hbar\,\delta_{ij}.
\end{align}
تتبادل الإحداثيات مع مركبات الزخم العائدة إلى اتجاهات مختلفة. فمثلًا، يتبادل \(x\) مع \(p_y\)، في حين إن \(x\) لا يتبادل مع \(p_x\).

احسب الآن
\begin{equation}
[L_x,L_y]
=
[yp_z-zp_y,\;zp_x-xp_z].
\end{equation}
تنعدم معظم الحدود. أما المساهمتان غير المنعدمتين فهما
\begin{align}
[yp_z,zp_x]&=-i\hbar\,yp_x,\\
[zp_y,xp_z]&=i\hbar\,xp_y.
\end{align}
لذلك
\begin{align}
[L_x,L_y]
&=i\hbar\left(xp_y-yp_x\right)\\
&=i\hbar L_z.
\end{align}
ويعطي التبديل الدوري لرموز الإحداثيات
\begin{align}
[L_x,L_y]&=i\hbar L_z,\\
[L_y,L_z]&=i\hbar L_x,\\
[L_z,L_x]&=i\hbar L_y.
\end{align}

\begin{classroomqa}
\lecturetimestamp{01:14:50}
\audiencequestion{عندما تعطي المبدلة \(L_z\)، هل لا يزال علينا أن نبيّن أن \(L_z\) ليس تركيبًا خطيًا من \(L_x\) و\(L_y\)؟}
\lecturerresponse{نعم. ثمة التزام بإثبات ذلك الاستقلال بدل افتراضه. ولا يُتابَع عند هذه النقطة برهان جبري مباشر.}
\end{classroomqa}

تنغلق المبدلات المعروضة على الفضاء المولد بـ \(L_x\) و\(L_y\) و\(L_z\): فمتى أُدرجت الثلاثة جميعًا، لم تعد المبدلات اللاحقة تنتج أي مؤثر خارج ذلك الفضاء. وهذا هو جبر لي للزخم الزاوي. ومع
\begin{equation}
[H,L_i]=0,
\qquad i=x,y,z,
\end{equation}
تتضمن هذه العلاقات قدرًا كبيرًا من المعلومات عن طيف الطاقة من دون الحاجة إلى حل معادلة Schr\"odinger الزاوية مباشرةً.

تُستخدم استراتيجية جبرية مشابهة في المذبذب التوافقي الكمومي، حيث تولد مؤثرات الخلق والفناء طيفه. وهنا ستغير المؤثرات المناظرة المركبة \(z\) من الزخم الزاوي بدلًا من الطاقة.

إن اختيار المحور \(z\) ليس إلا اختيارًا للتمثيل. وليس ثمة ما يجعل \(z\) مفضلًا فيزيائيًا على \(x\) أو \(y\)، لكن اختيار مركبة واحدة يعطينا أساسًا ملائمًا. اكتب
\begin{equation}
L_z\ket{m}=m\hbar\ket{m}.
\end{equation}
والمطلوب هو تحديد القيم التي يمكن أن يأخذها \(m\) وكيف ترتبط الحالات ذات القيم المختلفة لـ \(m\) بعضها ببعض.

\begin{classroomqa}
\lecturetimestamp{01:24:02}
\audiencequestion{عندما تقارن هذا البناء بالمذبذب التوافقي، هل تقصد المذبذب التوافقي الكمومي؟ فهذا لم يُغطَّ في المحاضرة السابقة.}
\lecturerresponse{نعم، المقارنة مع المذبذب الكمومي. وبرهان الزخم الزاوي الحالي قائم بذاته ولا يتطلب معرفة بالمذبذب. ويمكن لاحقًا قراءة المشابهة في الاتجاه الآخر: فالمذبذب يستخدم طريقة جبرية مشابهة.}
\end{classroomqa}

\section{رفع الزخم الزاوي وخفضه}

عرّف
\begin{equation}
L_+=L_x+iL_y,
\qquad
L_-=L_x-iL_y.
\end{equation}
وتستلزم مبدلات الزخم الزاوي
\begin{equation}
[L_z,L_+]=\hbar L_+,
\qquad
[L_z,L_-]=-\hbar L_-.
\end{equation}

طبّق العلاقة الأولى على حالة ذاتية لـ \(L_z\):
\begin{align}
L_zL_+\ket{m}
&=\left(L_+L_z+[L_z,L_+]\right)\ket{m}\\
&=(m+1)\hbar\,L_+\ket{m}.
\end{align}
فما لم يكن \(L_+\ket{m}=0\)، فإن الناتج حالة ذاتية أخرى لـ \(L_z\)، ولكن بقيمة ذاتية \((m+1)\hbar\). وبالمثل،
\begin{equation}
L_zL_-\ket{m}
=(m-1)\hbar\,L_-\ket{m}.
\end{equation}
إذن فإن \(L_+\) يرفع \(m\) بمقدار واحد و\(L_-\) يخفضه بمقدار واحد.

ابتداءً من أي قيمة ذاتية واحدة، يولد الجبر سلمًا تختلف قيمه بمقادير صحيحة. وقد يستمر هذا السلم بلا نهاية، أو قد ينتهي لأن
\begin{equation}
L_+\ket{m_{\mathrm{top}}}=0
\end{equation}
عند الطرف العلوي و
\begin{equation}
L_-\ket{m_{\mathrm{bottom}}}=0
\end{equation}
عند الطرف السفلي.

\begin{classroomqa}
\lecturetimestamp{01:35:11}
\audiencequestion{ألم يكن الزخم الزاوي قد فُرض بالفعل أن يكون عددًا صحيحًا بحجة الدائرة؟}
\lecturerresponse{تثبت تلك الحجة القيم الصحيحة للزخم الزاوي المداري الممثل بدوال موجية موضعية أحادية القيمة. أما جبر الزخم الزاوي المجرد فيعطي نتيجة أقل تقييدًا: فقد يحتوي السلم إما على قيم صحيحة أو على قيم نصف صحيحة.}
\end{classroomqa}

إذا انتهى سلم ما، فإن الثبات الدوراني يربط بين طرفيه. فدوران بمقدار \(180^\circ\) حول محور عمودي على \(z\) يحول
\begin{equation}
L_z\longrightarrow-L_z.
\end{equation}
ولذلك يجب أن يكون الطيف متناظرًا حول الصفر:
\begin{equation}
m_{\mathrm{bottom}}=-m_{\mathrm{top}}.
\end{equation}
وبما أن القيم المتجاورة تختلف بمقدار واحد، فثمة احتمالان. فإما أن يحتوي السلم على أعداد صحيحة،
\begin{equation}
\ldots,-2,-1,0,1,2,\ldots,
\end{equation}
أو على أنصاف أعداد صحيحة،
\begin{equation}
\ldots,-\frac{5}{2},-\frac{3}{2},-\frac{1}{2},
\frac{1}{2},\frac{3}{2},\frac{5}{2},\ldots,
\end{equation}
وتكون النهايتان حيثما أعطى مؤثر الرفع أو الخفض صفرًا.

وتنتقي حجة أحادية قيمة الدالة الموجية على الدائرة القيم الصحيحة للزخم الزاوي المداري،
\begin{equation}
\mathbf L=\mathbf r\times\mathbf p.
\end{equation}
أما القيم النصف الصحيحة فهي أيضًا فيزيائية، وتخص الزخم الزاوي المغزلي، الذي سيُعالَج على حدة.

وأخيرًا، تكون حالات السلم متساوية في الطاقة عندما يكون الهاملتوني ثابتًا تحت الدوران. افترض
\begin{equation}
H\ket{m}=E\ket{m}.
\end{equation}
وبما أن كل مركبة \(L_i\) تتبادل مع \(H\)، فإن \(L_+\) و\(L_-\) يتبادلان معه أيضًا. ومن ثم
\begin{align}
H L_+\ket{m}
&=L_+H\ket{m}\\
&=E L_+\ket{m}.
\end{align}
وكلما كان \(L_+\ket{m}\neq0\)، كانت للحالة المرفوعة الطاقة نفسها التي للحالة الأصلية. ويعطي الحساب نفسه
\begin{equation}
H L_-\ket{m}=E L_-\ket{m}.
\end{equation}
ولذلك يولد الرفع والخفض عائلة من الحالات ذات قيم ذاتية مختلفة لـ \(L_z\) ولكن بالطاقة نفسها تمامًا. وهذه هي الصلة بين التناظرات غير المتبادلة والانحلال التي جرى تطويرها على امتداد المحاضرة.

\begin{classroomqa}
\lecturetimestamp{01:44:13}
\audiencequestion{هل ترتبط حالات السلم هذه باختيار محاور دوران مختلفة؟}
\lecturerresponse{إنها مترابطة بالدورانات. فالمتجهات الذاتية لـ \(L_x\) أو \(L_y\) هي تراكيب خطية من المتجهات الذاتية المختارة لـ \(L_z\). ويغير تدوير حالة ما مركبات زخمها الزاوي مع إبقاء طاقتها من دون تغيير.}
\end{classroomqa}

إن اختيار \(L_+\) و \(L_-\) ملائم لأن هذين التركيبين يكشفان السلم مباشرةً. وتتكرر هذه الحيلة في ميكانيكا الكم كلها: فعندما يكون جبر المبدلات منغلقًا، فإن مؤثري الرفع والخفض ينظمان غالبًا طيف القيم الذاتية المرتبط به.

\begin{classroomqa}
\lecturetimestamp{01:46:33}
\audiencequestion{ماذا يحدث عندما يؤثر \(L_+\) في الحالة العظمى، وكيف يُعرَّف الطرف العلوي؟}
\lecturerresponse{إنه يعطي صفراً. وهذا هو تعريفه في سُلَّم منتهٍ. وبالمثل، يعطي \(L_-\) صفراً عند الطرف السفلي. ولا يلزم أن يكون السلم المجرَّد منتهيًا. أما حالات الزخم الزاوي الذرية التي سننظر فيها لاحقًا، فعادةً ما نصادف سلالم منتهية.}
\end{classroomqa}

الخطوة التالية هي المضي في تطوير بنية الزخم الزاوي هذه وتطبيقها على الذرات، بما في ذلك الهيدروجين. والمتعددات الدورانية التي وُجدت هنا تفسّر جزءًا من نمط الانحلال؛ ولا تزال ثمة بنية إضافية يتعين تطويرها.
